\documentclass[10pt,a4paper]{amsart}
\usepackage{amsmath,amssymb,amsthm,graphicx}
\usepackage{hyperref}
\usepackage{tikz}
%\usepackage{xcolor}
\usepackage{pict2e}
\let\arc\relax
%\usepackage{harpoon}
\usepackage{curves}
%\usepackage{chngcntr}
%\usepackage{apptools}
\usepackage{enumerate}

%\usepackage{verbatim}
%\usepackage[nospace,noadjust]{cite}
\usepackage{accents}
\usetikzlibrary{calc,decorations.markings}
\unitlength1cm

%%%%%%%%%%%%%%FONTS%%%%%%%%%%%%%%%%%%%%%%%%%%%%%%
\newcommand{\R}{\mathbb{R}}
\newcommand{\Z}{\mathbb{Z}}
\newcommand{\N}{\mathbb{N}}
\newcommand{\C}{\mathbb{C}}
\newcommand{\Id}{I}

%%%%%%%%%%%%%%MATHCAL%%%%%%%%%%%%%%%%%%%%%%%
\newcommand{\cK}{\mathcal{K}}
\newcommand{\cE}{\mathcal{E}}
\newcommand{\cR}{\mathcal{R}}
\newcommand{\cW}{\mathcal{W}}

%%%%%%%%%%%%%%%%%%ABBRS%%%%%%%%%%%%%%%%%%%%%%%%%%%%%

%%%%%%% OBJECTS SPECIFIC TO THIS PAPER%%%%%%%%%

\newcommand{\hw}{\widehat w}
\newcommand{\tw}{\widetilde w}

\newcommand{\bQ}{ Q}
\newcommand{\hQ}{{\widehat Q}}
\newcommand{\tQ}{{\widetilde Q}}

\newcommand{\sbQ}{Q^\star}
\newcommand{\shQ}{{\widehat Q^\star}}
\newcommand{\stQ}{{\widetilde Q^\star}}

\newcommand{\smu}{\mu^\star}
\newcommand{\shmu}{\widehat \mu^\star}
\newcommand{\stmu}{\widetilde \mu^\star}

\newcommand{\bmu}{{\mu}}
\newcommand{\hmu}{{\widehat \mu}}
\newcommand{\tmu}{{\widetilde \mu}}


\newcommand{\bP}{P}
\newcommand{\hP}{{\widehat P}}
\newcommand{\tP}{{\widetilde P}}

\newcommand{\sbP}{P^\star}
\newcommand{\shP}{{\widehat P^\star}}
\newcommand{\stP}{{\widetilde P^\star}}

\newcommand{\hR}{{\widehat R}}

\newcommand{\bv}{v}
\newcommand{\hv}{{\widehat v}}
\newcommand{\tv}{{\widetilde v}}

\newcommand{\sbv}{v^\star}
\newcommand{\shv}{{\widehat v^\star}}
\newcommand{\stv}{{\widetilde v^\star}}

\newcommand{\stu}{{\widetilde u^\star}}

\newcommand{\hF}{\widehat F}
\newcommand{\hW}{\widehat W}

\newcommand{\Fw}{\frac{F^2}{w}}
\newcommand{\hFw}{\frac{\hF^2}{\hw}}

\newcommand{\Cau}{\mathcal C_{\Sigma}}
\newcommand{\Cauv}{\mathcal C_{v}}
\newcommand{\Cauhv}{\mathcal C_{\hv}}
\newcommand{\Cautv}{\mathcal C_{\tv}}

\newcommand{\sCau}{\mathcal C_{\Sigma^\star}}
\newcommand{\sCauv}{\mathcal C_{\bv^\star}}
\newcommand{\sCauhv}{\mathcal C_{\hv^\star}}
\newcommand{\sCautv}{\mathcal C_{\stv}}

\newcommand{\Ltwo}{L^2(\Sigma)}
\newcommand{\sLtwo}{L^2(\Sigma^\star)}

\newcommand{\Sigl}{\Sigma^\ell}
\newcommand{\Sigr}{\Sigma^r}

%%%%%%%%%%%%%%%%OTHER SHORTCUTS%%%%%%%%%%%
\newcommand{\nn}{\nonumber}
\newcommand{\ul}{\underline}
\newcommand{\ol}{\overline}
\newcommand{\pa}{\partial}
\newcommand{\ti}{\tilde}
\newcommand{\wti}{\widetilde}
\newcommand{\wha}{\widehat}
\newcommand{\norm}[1]{\lVert#1 \rVert}
\newcommand{\abs}[1]{\lvert#1 \rvert}
%\newcommand{\Id}{\mathbb{I}}
\newcommand{\I}{\mathrm{i}}
\newcommand{\E}{\mathrm{e}}
\newcommand{\re}{\mathop{\mathrm{Re}}}
\newcommand{\im}{\mathop{\mathrm{Im}}}
\newcommand{\eps}{\varepsilon}
\newcommand{\dist}{\mathop{\mathrm{dist}}}

\newtheorem{Theorem}{Theorem}[section]
\newtheorem{Corollary}[Theorem]{Corollary}
\newtheorem{Lemma}[Theorem]{Lemma}
\newtheorem{Proposition}[Theorem]{Proposition}
 { \theoremstyle{definition}
\newtheorem{Definition}[Theorem]{Definition}
\newtheorem{Note}[Theorem]{Note}
\newtheorem{Example}[Theorem]{Example}
\newtheorem{Remark}[Theorem]{Remark} }

\numberwithin{equation}{section}
\usepackage{cleveref}
\crefname{originalequation}{original equation}{original equations}
\Crefname{originalequation}{Original equation}{Original equations}
\makeatletter
\newlabel{3/2decay}{{internal}{source}{}{sourcepointer-1}{}}
\newlabel{3/2decay@cref}{{[originalequation][0][]internal}{[1][0][]source}{}{}{}}
\makeatother
\crefname{originalequation}{original equation}{original equations}
\Crefname{originalequation}{Original equation}{Original equations}
\makeatletter
\newlabel{ParaBehaviour}{{internal}{source}{}{sourcepointer-2}{}}
\newlabel{ParaBehaviour@cref}{{[originalequation][0][]internal}{[1][0][]source}{}{}{}}
\makeatother
\crefname{originalequation}{original equation}{original equations}
\Crefname{originalequation}{Original equation}{Original equations}
\makeatletter
\newlabel{ParaConv}{{internal}{source}{}{sourcepointer-3}{}}
\newlabel{ParaConv@cref}{{[originalequation][0][]internal}{[1][0][]source}{}{}{}}
\makeatother
\crefname{originalequation}{original equation}{original equations}
\Crefname{originalequation}{Original equation}{Original equations}
\makeatletter
\newlabel{phiEndpoints}{{internal}{source}{}{sourcepointer-4}{}}
\newlabel{phiEndpoints@cref}{{[originalequation][0][]internal}{[1][0][]source}{}{}{}}
\makeatother
\title{The integral of the singular part of the jump matrix is of cubic order}
\author{Deift, Percy, Piorkowski, Mateusz}
\date{doi:10.3842/SIGMA.2024.004; CC BY 4.0}
\begin{document}
\maketitle
\setcounter{section}{7}
\setcounter{Theorem}{5}
\noindent\textit{Editorial note.} All delivered mathematical text is the author's own, copied from the pinned SIGMA article source; no statement or proof step was rewritten, reordered or removed. The counted statement is the labeled lemma environment at source lines 2008-2013 and the counted proof is the proof environment at source lines 2014-2049.\par\medskip
\noindent\textit{Editorial note.} The proof cites published references of the author's bibliography by key; the cited entries are copied line for line from the author's own bibliography into a retained bibliography fragment, and hyperref is added because the publisher class, not the author preamble, supplied the link command.\par\medskip
\noindent\textit{Editorial note.} The retained preamble of this unit was generated mechanically for this delivery: it copies the author's own package, numbering, macro and theorem declaration block from the pinned source (source lines 9-133), with the publisher class and page setup lines removed.\par\medskip
\noindent\textit{Allowed standard prerequisites (editorial).}\begin{itemize}
\item The matching relation between the parametrices and the outer parametrix. same article, the display at source line 1377. Bounds the quotients appearing in the integrand and their inverses uniformly.
\item The growth bounds of the local parametrices near the endpoint. same article, the display at source line 1381 and the decay estimate at source line 541. Provide the inverse square root bound and the exponential decay that together make the integrand small.
\item The behaviour of the phase function at the endpoints. same article, the display at source line 429. Converts the distance to the endpoint into the variable in which the exponential decay is expressed.
\end{itemize}\medskip
\ifdefined\hypertarget\hypertarget{sourcepointer-1}{}\fi
\noindent\textit{Source reference.} Original equation internal: sources/sigma-2024-004-recurrence-coefficients-for-orthogonal-polyn/pinned-original-source.tex, source lines 541--545. The counted statement or proof refers to that the author's displayed decay estimate for the exponential factor; it is the author's own earlier display, definition or statement in the same article and the whole argument of the counted span is retained verbatim.\par\medskip
\ifdefined\hypertarget\hypertarget{sourcepointer-2}{}\fi
\noindent\textit{Source reference.} Original equation internal: sources/sigma-2024-004-recurrence-coefficients-for-orthogonal-polyn/pinned-original-source.tex, source lines 1381--1383. The counted statement or proof refers to that the author's displayed growth bound for the local parametrices near the endpoint; it is the author's own earlier display, definition or statement in the same article and the whole argument of the counted span is retained verbatim.\par\medskip
\ifdefined\hypertarget\hypertarget{sourcepointer-3}{}\fi
\noindent\textit{Source reference.} Original equation internal: sources/sigma-2024-004-recurrence-coefficients-for-orthogonal-polyn/pinned-original-source.tex, source lines 1377--1379. The counted statement or proof refers to that the author's displayed matching relation between the parametrices and the outer parametrix; it is the author's own earlier display, definition or statement in the same article and the whole argument of the counted span is retained verbatim.\par\medskip
\ifdefined\hypertarget\hypertarget{sourcepointer-4}{}\fi
\noindent\textit{Source reference.} Original equation internal: sources/sigma-2024-004-recurrence-coefficients-for-orthogonal-polyn/pinned-original-source.tex, source lines 429--433. The counted statement or proof refers to that the author's displayed behaviour of the phase function at the endpoints; it is the author's own earlier display, definition or statement in the same article and the whole argument of the counted span is retained verbatim.\par\medskip
\begin{Lemma}\label{FinalLemma}
For $n \to \infty$, we have
\begin{align}\label{FinalInt}
 \int_{ \Sigma^U} \bmu_n(\bv_n-\hv_n)\hmu_n^{-1} {\rm d}s = O\left(\frac{1}{n^3}\right).
\end{align}
\end{Lemma}

\begin{proof}
First, we need to show that $\sbQ(z) = \bQ(z)[\bP(z)]^{-1}$ for $z \in U$ is uniformly bounded as $n \to \infty$. Note that this statement is more intricate than the corresponding statements for $\hQ(z)\big[\hP(z)\big]^{-1}$ and $\tQ(z)\big[\tP(z)\big]^{-1}$, as the logarithmic weight function is not covered in \cite{KMVV} and hence we do not have equation~\eqref{ParaConv} for $\bQ(z)[\bP(z)]^{-1}$ at our disposal. Recall that for~${s\!\in\! \Sigma^\ell = \partial U}$, we have
\begin{align}\label{smu_on_U}
 \smu(s) = \sbQ_-(s) = \bQ(s)\big[\bP(s)\big]^{-1}.
\end{align}
From \eqref{smu_on_U}, we see that on $\partial U$, $\smu$ is in fact the restriction of the analytic function $\sbQ(z) = \bQ(z)[\bP(z)]^{-1}$, $z\in U$. Moreover, $\Vert \smu \Vert_{L^2( \Sigma^U)} = O(1)$ for $n \to \infty$, as $\smu = (1-\sCauv)^{-1}\Id$ and~${(1-\sCauv)^{-1}}$ is uniformly bounded on $L^2(\Sigma^\star)$. Thus, using Cauchy's integral formula, we get for~$z \in U$
\begin{align*}
 |\sbQ(z)| = \bigg|\frac{1}{2\pi\I} \int_{\partial U} \frac{\smu(s){\rm d}s}{s-z}\bigg| \lesssim O\left(\frac{1}{\dist(z,\partial U)}\right).
\end{align*}
We have some freedom to choose $U$, so if necessary we can shrink it and conclude
\begin{align}\label{Qstar}
 |\sbQ(z)| = O(1), \qquad z\in U,
\end{align}
uniformly as $n \to \infty$. Let us now rewrite \eqref{FinalInt}, where we drop the $n$-dependence for better readability:
\begin{align}\label{FinalInt2}
 \int_{ \Sigma^U} \bmu(\bv-\hv)\hmu^{-1} {\rm d}s = \int_{ \Sigma^U} \sbQ(s)\bP_-(s)(\bv(s)-\hv(s))\big[\hP_-(s)\big]^{-1}\big[\shQ(s)\big]^{-1} {\rm d}s.
\end{align}
We now list bounds for each of the factors in the integrand:
\begin{itemize}\itemsep=0pt
 \item \eqref{phiEndpoints} and \eqref{3/2decay} imply that for some $c > 0$
 \begin{align*}\bv(s)-\hv(s) =\begin{cases} \begin{pmatrix}
 0 & 0
 \\
 \big(\Fw(s)-\hFw(s)\big)\phi^{-2n}(s) & 0
 \end{pmatrix}
 \\
 \quad= O\big(|s+1|^{3/2}\E^{-cn|s+1|^{1/2}}\big),\qquad s \in \Sigma^U \setminus (-1,1),
 \\
  0, \qquad s \in \Sigma^U \cap (-1,1),
 \end{cases}
 \end{align*}
 \item \eqref{ParaConv} and \eqref{Qstar} implies that $[\sbQ(z)]^{\pm 1}, [\shQ(z)]^{\pm 1} = O(1)$,
 \item and \eqref{ParaBehaviour} implies that $\big[\bP_-(z)\big]^{\pm 1}, \big[\hP_-(z)\big]^{\pm 1} = O\big(|z+1|^{-1/2}\big)$.
\end{itemize}
Taking the contributions of all factors in \eqref{FinalInt2} into account we see that the integrand can be estimated by $O\big(|z+1|^{1/2}\E^{-cn|z+1|^{1/2}}\big)$, which precisely integrates to the error in \eqref{FinalInt}.
\end{proof}

\begin{thebibliography}{99}
\bibitem{KMVV}
Kuijlaars A.B.J., McLaughlin K.T.-R., Van~Assche W., Vanlessen M., The
 {R}iemann--{H}ilbert approach to strong asymptotics for orthogonal
 polynomials on {$[-1,1]$}, \href{https://doi.org/10.1016/j.aim.2003.08.015}{\textit{Adv. Math.}} \textbf{188} (2004), 337--398,
 \href{https://arxiv.org/abs/math.CA/0111252}{arXiv:math.CA/0111252}.
\end{thebibliography}
\end{document}
